\documentclass{beamer}

\usepackage{xcolor}
\usepackage{tikz}

% Theme choice
\usetheme{Boadilla}
% 120,180,250
\definecolor{MyBlue}{RGB}{0,40,255}
\usecolortheme[named=MyBlue]{structure}

% Packages
\usepackage[utf8]{inputenc}
\usepackage{graphicx}

% Title page info
\title{TODO}
\author{Noah Albers, Felix Graff}
\institute{Uni Münster}
\date{\today}

\usepackage{comment}
\newcommand{\teacher}[1]{\textcolor{red}{#1}}
\begin{document}
	
	% Title slide
	\frame{\titlepage}
	
	
	\begin{frame}{Übersicht}
		\tableofcontents
	\end{frame}
	
	
	% Section example
	\section{Differenzierung - Was ist das?}
	
	\begin{frame}[t]{Differenzierung - Was ist das?}{Äußere Differenzierung}
		\begin{itemize}
			\item Teilung in möglichst \textit{homogene} Gruppen
			\item Basiert oft auf:
			\begin{itemize}
				\item Leistung (Noten)
				\item Interessen (Fächer)
			\end{itemize}
			\teacher{HIER FRAGEN: Fallen euch Beispiele für äußere Diff. ein?}
			\pause
			
			\item Mögliche Beispiele Äußerer Differenzierung
			\begin{itemize}
				\item Schulform entscheided – Hauptschule, Realschule, Gymnasium
				\item Fremdsprache
				\item Fachbelegung in der Oberstufe
			\end{itemize}
			$\Rightarrow$ Kann von der Lehrkraft kaum beeinflusst werden
		\end{itemize}



	\end{frame}


	
	\begin{frame}[t]{Differenzierung - Was ist das?}{Innere Differenzierung}
		\begin{itemize}
			\item Bezieht sich auf \textit{heterogene} Gruppen
			\item Maßnahmen die in den Gruppen getroffen werden 
			\begin{itemize}
				\item Auf Fähigkeiten der SuS eingehen
				\item Fertigkeiten der SuS eingehen 
				\item Den Interessen und Bedürfnisse entsprechen
				\item Individuelle Lernbedürfnisse und Wünsche beachten
			\end{itemize}
			\item Äußere Differenzierung reicht in der Regel nicht (kaum Umsetzbar)\\$\Rightarrow$ es braucht Innere Differenzierung
		\end{itemize}
	\end{frame}
	
	
	
\end{document}